\documentclass[12pt,a4paper]{article}
\usepackage{../../../template/stile}

%Titolo documento
\newcommand{\titoloDocumento}{Norme di Progetto}

%Prima data di creazione del documento
\newcommand{\dataCreazione}{30 novembre 2015}

%Inserite la versione attuale del documento
\newcommand{\versione}{1.0.3}

%Stato in cui si trova il documento: Formale solo all'atto di consegna
\newcommand{\stato}{Informale}

%Uso del documento
\newcommand{\uso}{Interno}

\rhead{\titoloDocumento}
\lfoot{Versione: \versione}
\title{\titoloDocumento}

\begin{document}
\begin{titlepage}
\begin{center}
\today \\
\vspace{1cm}
\begin{Huge}
\textbf{\nomeGruppo} \\
\end{Huge}
\textbf{\prjL} \\
\vspace{1cm}
\includegraphics[scale=0.5]{\logoGrande}
\vspace{1cm}

\HRule \\[0.4cm]
\begin{Huge}
{\huge \bfseries \titoloDocumento}\\[0.4cm]
\end{Huge}
\HRule \\[1cm]
\vfill

\begin{table}[h]
\begin{center}
\begin{tabular}{r | l}
\multicolumn{2}{c}{\textbf{Informazioni sul documento}}\\
\midrule
\textbf{Nome Documento}	&	\titoloDocumento	\\
\textbf{Versione}	&	\versione	\\
\textbf{Stato}	&	\emph{\stato}	\\
\textbf{Uso}	&	\emph{\uso}	\\
\textbf{Data Creazione}	&	\dataCreazione	\\
\textbf{Data Ultima Modifica}	&	\today	\\
\textbf{Redazione}	&	\NDC	\\
\ &	\AVE	\\
\ &	\AVI	\\
\textbf{Verifica}	&	\IB	\\
\ & \TP \\
\textbf{Approvazione}	&	\AB	\\
\textbf{Lista Distribuzione}	&	\nomeGruppo	\\

\end{tabular}
\end{center}
\end{table}

\end{center}
\end{titlepage}
\newpage

\Large{\textbf{Registro delle modifiche}}
\normalsize

\begin{table}[h]
\begin{center}

\begin{tabular}{p{0.12\textwidth} p{0.2\textwidth} p{0.18\textwidth} p{0.39\textwidth}}
\toprule
\textbf{Versione}	&	\textbf{Autore}	&	\textbf{Data}	&	\textbf{Descrizione}\\
\midrule
\midrule
1.0.4 & \AVE & 2015-12-08 & Ruoli di progetto e Riunioni \\
\midrule
1.0.3 & \AVI & 2015-12-07 & Norme del Piano di Progetto e dell'Analisi dei requisiti \\
\midrule
1.0.2 & \AVI & 2015-12-03 & Descrizione processi di base \\
\midrule 
1.0.1 & \NDC & 2015-12-02 & Introduzione \\
\midrule
1.0.0 & \NDC & 2015-11-30 & Template di base \\
\bottomrule
\end{tabular}
\caption{Versionamento del documento}
\label{tabVers1}
\end{center}
\end{table}
\newpage

\tableofcontents
\newpage

\listoftables
\listoffigures
\newpage

% NOTA non so dove mettere l'uso di git

\section{Introduzione}
Questo documento definisce le norme che i membri del gruppo \nomeGruppo{} adotteranno nello svolgimento del progetto \prjL. Tutti i membri sono tenuti a leggere il documento attentamente e a seguirene le norme.

\subsection{Scopo}
Il documento si propone come proposito di garantire l’uniformità di tutto il materiale prodotto, migliorarene l’efficienza e ridurrene gli errori.

\subsection{Descrizione}
Verranno definite norme riguardanti:
\begin{itemize}
  \item interazioni tra membri del gruppo;
  \item stesura documenti, convenzioni stilistiche e tipografiche;
  \item modalità di lavoro durante tutte le fasi del progetto;
  \item ambiente di lavoro.
\end{itemize}

\subsection{Glossario}
Al fine di evitare ogni ambiguità di linguaggio e massimizzare la comprensione dei
documenti, i termini tecnici, di dominio, gli acronimi e le parole che necessitano chiarimenti sono riportate nel documento \emph{Glossario\_v1.2.0}. Ogni occorrenza di vocaboli presenti nel \GlO è marcata da una G maiuscola in pedice (e.g. \gloss{pedice}).

\newpage

\section{Processi base} %ISO 12207 5
I processi base sono formati dai seguenti sotto processi:
\begin{enumerate}
\item Processo dell’acquirente
\item Processo del fornitore
\item Processo di sviluppo
\item Processo operativo
\item Processo di manutenzione
\end{enumerate}
Per lo scopo del corrente progetto saranno dettate solo le norme per le fasi del fornitore e di sviluppo, poiché i processi  dell’acquirente non riguardano il team e le fasi operative e di manutenzione non verranno trattate.

\subsection{Processi del fornitore}

\subsubsection{Scopo}
Definisce le attività e i task attraverso i quali il fornitore comunica al committente modo, tempi e costi previsti per il progetto.

\subsubsection{Descrizione}
Consiste nel redigere il \PdQ, che dovrà illustrare la strategia complessiva di verifica e validazione proposta dal fornitore per pervenire al collaudo del sistema con la massima efficienza ed efficacia, e il \PdP, che dovrà presentare l'organigramma dettagliato del fornitore, lo schema proposto per l'assegnazione e la rotazione dei ruoli di progetto, l'impegno complessivo previsto per ogni ruolo e per ogni individuo, e il conto economico preventivo di realizzazione del prodotto. Le regole dettagliate per redigere il \PdP\ verranno analizzate nel prossimo paragrafo.

\subsubsection{Piano di Progetto} %ISO 12207 5.2.4
Il \PM, assieme all'\AM, deve presentare nel \PdP\ diagrammi di Gant e ripartizione delle risorse relativi ad ogni fase del Progetto. In Particolare è necessario che ogni componente del gruppo ricopra tutti i ruoli nell'arco dello svolgimento del progetto. La rotazione dei ruoli deve garantire un'equa ripartizione del carico di lavoro individuale, ovvero il totale di ore produttive per ogni persona può differire al più di poche unità da quello degli altri. Inoltre è possibile che un componente rivesta più ruoli contemporaneamente, ma sempre evitando conflitti di interesse. In particolare i conflitti di interesse da evitare sono quelli tra \PM\ e qualsiasi altro ruolo o tra \VR\ e qualsiai altro ruolo.

\subsection{Processi di sviluppo}

\subsubsection{Scopo}
Il processo di sviluppo produce un prodotto software che soddisfi i requisiti architetturali. Inoltre definisce un insieme di azioni che formalizzano comportamenti, interfacce e vincoli di implementazione atti a creare un prodotto software

\subsubsection{Descrizione}
Il processo di sviluppo è un insieme di atività e task necessari per lo sviluppo di un prodotto software. Coloro che eseguono queste attivita’ sono gli sviluppatori. Le attività che compongono i processi di sviluppo sono:
\begin{enumerate}
\item Processo di Analisi dei requisiti
\item Processo di Progettazione
\item Processo di Codifica
\end{enumerate}

\subsubsection{Analisi dei requisiti} %ISO 12207 5.3.2

\paragraph{Scopo}
Gli Analisti dovranno produrre l’Analisi dei Requisiti basandosi sul capitolato e sugli incontri con il proponente. Tale processo ha l’obiettivo di formalizzare e rendere tracciabile in un documento i requisiti e casi d’uso individuati, comprendendo a fondo eventuali problemi da risolvere in fase di progettazione.

\paragraph{Descrizione}
Con il completamento di questo processo si ottiene una documentazione affidabile e consistente che ben descrive le esigenze e le richieste del proponente.

\paragraph{Classificazione dei Requisiti}
È compito degli Analisti stilare una lista dei requisiti emersi dal capitolato e da eventuali
riunioni con il proponente. Questi requisiti dovranno essere classificati per tipo e
importanza, utilizzando la seguente codifica:
\begin{center}
R[importanza/tipo/codice]
\end{center}
\begin{itemize}
\item Importanza può assumere i seguenti valori:
\begin{description}
\item[0] Requisito obbligatorio
\item[1] Requisito desiderabile
\item[2] Requisito opzionale
\end{description}
\end{itemize}
\begin{itemize}
\item Tipo può assumere i seguenti valori:
\begin{description}
\item[F] Funzionale
\item[Q] di Qualità
\item[P] Prestazionale
\item[V] Vincolo
\end{description}
\end{itemize}
\begin{itemize}
\item Codice è un numero progressivo univoco per ogni requisito, indipendente da importanza e tipo. Nel caso si abbia un sottorequisito codice può anche essere espresso in modo gerarchico tramite la notazione \textit{CodiceRequistoPadre.CodiceSottorequisito}
\end{itemize}
Ogni requisito deve essere correlato da una sintetica ma precisa descrizione

\paragraph{Casi d'uso}
I casi d'uso sono creati secondo lo standard UML 2.0 e sono identificati dalla seguente notazione:
\begin{center}
UC[codice]
\end{center}
dove il codice è un nuero progressivo identificativo di ogni requitito, gerarchico nel  caso di sottocasi d'uso tramite la notazione \textit{CodiceUCPadre.CodiceSottoUC}. Per ogni caso d'uso deve inoltre essere indicato:
\begin{itemize}
\item \textbf{Titolo}
\item \textbf{Attori} principali e secondari coinvolti
\item \textbf{Descrizione} sintetica del Caso d'uso
\item \textbf{Precondizione}
\item \textbf{Postcondizione}
\item \textbf{Requisiti} collegati al Caso d'uso
\end{itemize}
Per la creazione di diagrammi UML verrà usato il software ArgoUML (\textit{argouml.tigris.com})

\subsubsection{Progettazione}
Le norme di progetto saranno ampliate con questa sezione non appena verranno concordate dal gruppo.

\subsubsection{Codifica}
Le norme di progetto saranno ampliate con questa sezione non appena verranno concordate dal gruppo.

\newpage

\section{Processi di Supporto} %ISO 12207 6

\subsection{Scopo}
Definisce norme per lo sviluppo e il mantenimento della documentazione prodotta durante il ciclo di vita del software. Inoltre definisce metodi per il controllo della qualità, di verifica e validazione di tali documenti.

\subsection{Descrizione}
Verranno definite norme riguardanti:
\begin{itemize}
  \item utilizzo e l'accesso ai documenti;
  \item templates da utilizzare, l'uniformità di linguaggio, le convenzioni stilistiche e tipografiche;
  \item metodi di verifica e approvazione dei documenti;
\end{itemize}

\subsection{Documentazione}\label{Documentazione} %ISO 12207 6.1
La readazione di tutti i documenti viene assegnata dal \PM{} ai responsabili della stesura di tale documento secondo i \gloss{ruoli} definiti in \TODO{definire i ruoli}. Successivamente se i redattori ritengono che il documento sia completo, su comferma da parte del \PM, il documento deve essere preso in visione da un \VR. Se il \VR ritiene che documento rispetta i requisiti imposti dal \PR{} e supera il controllo di qualità può proporre al \PM{} l'approvazione di tale documento. Se il documento viene approvato è giunto alla fase finale.

Tutti documenti essi dovranno a essere denominati secondo il seguente formalismo:
\begin{center}
\emph{NomeDocumento\_vX.Y.Z.pdf}
\end{center}
dove \emph{NomeDocumento} rappresenta il nome del documento e \emph{vX.Y.Z} la versione come spiegata nella sezione \ref{Versione}. Il documento non può contentre lettere accentate. Nel caso il nome fosse composto da più parole, la prima lettera di ogni parola deve essere maiuscola e non devono esserci spazi, underscore o trattini a separare le parole.

Ogniqualvolta sia necessaria la citazione di una versione specifica di un documento, essa deve comprendere sia il nome che il numero di versione aderente al formato:
\begin{center}
\emph{Nome Documento vX.Y.Z}
\end{center}

\subsubsection{Versione} \label{Versione}
La documentazione prodotta deve essere corredata del numero di versione attuale tramite la seguente codifica \emph{vX.Y.Z} dove:
\begin{itemize}
  \item \textbf{X} indica il numero crescente di uscite formali del documento. All’inizio di ogni fase il \PM deve cambiare tale indice seguendo la numerazione progressiva indicata e impostare a 0 gli indici Y e Z. L’indice deve seguire la numerazione progressiva indicata e non sono ammessi indici diversi da quelli elencati:
  \begin{enumerate}
    \item fase che si conclude con la \RR;
    \item fase che si conclude con la \RP;
    \item fase che si conclude con la \RQ;
    \item fase che si conclude con la \RA;
  \end{enumerate}
  \item \textbf{Y} indica la fase in cui il documento si trova. Nel momento in cui inizia l’attività di stesura il redattore del documento deve controllare che tale indice sia correttamente impostato a 0. All’inizio della verifica il \VR{} deve variare l’indice impostandolo a 1, dopo aver ricevuto il consenso dal \PM. Conclusa la verifica, il \PM{} provvede all’approvazione del documento e deve impostare l’indice a 2. L’indice deve seguire la numerazione progressiva indicata e non sono ammessi indici diversi da quelli elencati:
  \begin{enumerate}[start=0]
    \item stesura del documento;
    \item verifica del documento;
    \item approvazione del documento.
  \end{enumerate}
  \item \textbf{Z} indica il numero crescente di modifiche minori. Ad ogni modifica effettuata deve corrispondere ad un’aggiunta di una voce nel diario delle modifiche, il redattore o il \VR devono aggiornare l’indice seguendo una numerazione progressiva. Non viene fissato un limite superiore per tale indice.
\end{itemize}

\subsubsection{Stato}
I documenti sono provvisti di uno \gloss{stato}. Tale stato può essere \gloss{informale} o  \gloss{formale}.

\subsubsubsection{Informale}
Tutti i documenti saranno da ritenersi informali fino all’approvazione del \PM, il quale potrà richiederne una revisione ulteriore. L’utilizzo dei documenti informali è da considerarsi esclusivamente interno al gruppo e localizzato durante la fase di redazione e verifica di tali documenti.

\subsubsubsection{Formale}
I documenti approvati dal \PM{} si riterranno formali e pronti per essere distribuiti. Solo i documenti formali potranno essere distribuiti alla loro lista di distribuzione. Ogniqualvolta un documento formale venga modificato o rivisitato, la nuova versione è da considerarsi non formale fino ad approvazione del \PM e quindi sarà trattata come un documento informale.

\subsubsection{Uso}
I documenti possono avere diverse liste di distribuzione ma una devono essere classificati in due principali categorie: uso interno ed esterno.

\subsubsubsection{Interno}
I docuemnti definiti ad uso interno non devono essere distribuiti all'esterno del gruppo di lavoro.

\subsubsubsection{Esterno}
I docuemnti definiti ad uso esterno devono essere distribuiti secondo i criteri della lista di distribuzine solamente se sono stati approvati dal \PM{} e se sono in stato \gloss{formale}.

\subsubsection{Struttura dei documenti}
I documenti devono rispettare una struttura prefissata. Per agevolare la redazione della documentazione è stato creato un template \LaTeX{} contenente tutte le impostazioni stilistiche e grafiche citate in questo documento. Tale modello si può trovare nella \gloss{repository} in \filePath{doc/template}.

\subsubsubsection{Prima pagina}
Ogni documento deve avere una prima pagina che contiene le seguenti informazioni sul documento:
\begin{itemize}
  \item nome del gruppo;
  \item nome del progetto;
  \item logo del gruppo;
  \item titolo del documento;
  \item versione del documento;
  \item stato del documento (\gloss{informale} o \gloss{formale})
  \item uso del documento (\gloss{interno} o \gloss{esterno});
  \item data di creazione del documento;
  \item data di dultima modifica del documento;
  \item nome e cognome dei redattori del documento;
  \item nome e cognome dei verificatori del documento;
  \item nome e cognome del responsabile che approva il documento;
  \item lista di distribuzione del documento;     
\end{itemize}

\subsubsubsection{Diario delle modifiche}
Il diario delle modifiche del documento segue la prima pagina. La tabella deve essere ordinata per data in ordine decrescente, in modo che la prima riga corrisponda alla versione attuale del documento. Ogni riga del diario delle modifiche deve contenere:
\begin{itemize}
  \item versione del documento dopo la modifica;
  \item nome e cognome dell'autore della modifica;
  \item data della modifica;
  \item titolo del documento;
  \item una breve descrizione delle modifiche svolte.
\end{itemize}

\subsubsubsection{Indice}
Al diario delle modifiche deve sempre seguire un indice delle sezioni e sottosezioni del documento.

\subsubsubsection{Elenco tabelle e figure}
All'indice delle sezioni e sottosezioni può seguire un elenco delle tabelle e delle figure. Nel caso non siano presenti figure o tabelle i rispettivi indici devono essere omessi.

\subsubsubsection{Introduzione}
Ogni documento deve essere provisto di un'introduzione che ne spiega brevemente il contenuto. L'introduzione deve contenere tue sottosezioni denominate \emph{Scopo} e \emph{Descrizione}. Lo scopo definisce a che proposito viene scritto il documento mentre la descrizione ne descrive brevemente gli argomenti trattati.

\subsubsubsection{Sezioni}
I documenti devono avere altre sezioni in modo da orgnanizzare meglio i contenuti. Ogni sezione deve avere tue sottosezioni denominate \emph{Scopo} e \emph{Descrizione}. Lo scopo definisce a che proposito viene scritta la sezione mentre la descrizione ne descrive brevemente gli argomenti trattati. Le sezioni possono inoltre avere altre sottosezioni. Queste sottosezioni possono avere al loro interno gerarchie annidate di sottosottosezioni,  paragrafi e sottoparagrafi. Questi non necessitano di \emph{Scopo} e \emph{Descrizione}.

\subsubsubsection{Formattazione delle pagine}
L’intestazione di ogni pagina, apparte la prima, deve contenere:
\begin{itemize}
  \item logo del gruppo;
  \item nome del documento.
\end{itemize}
Il piè di pagina deve contenere:
\begin{itemize}
  \item versione del documento;
  \item nome dell'università e anno accademico corrente;
  \item pagina corrente nel formato N di T dove N è il numero di pagina corrente e T è il numero di pagine totali;
  \item email del gruppo;
  \item licenza di distribuzione del documento.
\end{itemize}

\subsubsection{Norme tipografiche}

\subsubsubsection{Stile del testo}
\begin{itemize}
  \item \textbf{Corsivo}: deve essere usato solo per indicare termini in lingua inglese, standard, citazioni, \gloss{ruoli} e nomi di documenti
  \item \textbf{Grassetto}: deve essere usato solo per parole significative che devono essere messe in risalto. Il grassetto deve essere usato con estrema cautela poichè può appesantire la lettura. Va applicato ai titoli e agli elementi che costituiscono un elenco se sono seguiti da una descrizione
  \item \textbf{Monospace}: deve essere usato solo per indicare i percorsi e i nomi dei file, codice e comandi
  \item \textbf{Maiuscolo}: deve essere usato solo per gli acronimi
\end{itemize}

\subsubsubsection{Composizione}
\begin{itemize}
  \item \textbf{Elenchi}: ogni punto dell’elenco deve terminare con un punto e virgola, tranne l’ultimo che deve terminare con un punto. Gli elenchi devono essere puntati se non è necessario che gli elementi siano in un ordine preciso. Se l'elenco è una lista di elementi seguiti da descrizione devono terminare senza punteggiatura
  \item \textbf{Membri del gruppo}: i nomi dei membri del grupppo si devono indicare con nome e poi cognome. Sono state definite macro in \LaTeX{} con tutti i nomi del gruppo per facilitare la composizione dei documenti:
  \begin{itemize}
    \item \textbf{Alessio Vitella}: \latex{\textbackslash AVI}
    \item \textbf{Andrea Venier}: \latex{\textbackslash AVE}
    \item \textbf{Nicola De Cao}: \latex{\textbackslash NDC}
    \item \textbf{Igor Baylyak}: \latex{\textbackslash IB}
    \item \textbf{Walter Sandon}: \latex{\textbackslash WS}
    \item \textbf{Thomas Pigarelli}: \latex{\textbackslash TP}
    \item \textbf{Anna Bonaldo}: \latex{\textbackslash AB}
  \end{itemize}
  \item \textbf{Ruoli}: i nomi dei ruoli devono essere riportati in corsivo e con la prima lettera maiuscola. Sono state definite macro in \LaTeX{} con tutti i ruoli definti per facilitare la composizione dei documenti:
  \begin{itemize}
    \item \textbf{Amministratore}: \latex{\textbackslash AM}
    \item \textbf{Analista}: \latex{\textbackslash AN}
    \item \textbf{Progettista}: \latex{\textbackslash PG}
    \item \textbf{Programmatore}: \latex{\textbackslash PR}
    \item \textbf{Verificatore}: \latex{\textbackslash VR}
    \item \textbf{Project Manager}: \latex{\textbackslash PM}
  \end{itemize}
  \item \textbf{Glossario}: ogni termine che possa causare ambiguità deve essere marcato da una lettera G scritta a pedice e riportata nel \GlO. È stata fatta una definite macro in \LaTeX{} per facilitare la composizione dei documenti: \latex{\textbackslash gloss}
  \item \textbf{Note}: le note a piè di pagina possono contebnere link o URL e devono essere indicate con il comando \LaTeX: \latex{\textbackslash footnote}
  \item \textbf{\LaTeX}: i riferimenti a \LaTeX{} devono essere fatti utilizzato il comando \LaTeX: \latex{\textbackslash LaTeX}
  \item \textbf{Revisioni}: i riferimenti alle revisioni devono essere fatti utilizzato la rispettiva macro che garantisce la corretta scrittura con la prima lettera di ogni parola che non sia una preposizione maiuscola :
  \begin{itemize}
    \item \textbf{Revisione dei Requisiti}: \latex{\textbackslash RR}
    \item \textbf{Revisione di Progettazione}: \latex{\textbackslash RP}
    \item \textbf{Revisione di Qualifica}: \latex{\textbackslash RQ}
    \item \textbf{Revisione di Accettazione}: \latex{\textbackslash RA}
  \end{itemize}
  \item \textbf{Documenti}: i riferimenti ai documenti devono essere fatti utilizzato la specifica come definita nella sezione \ref{Documentazione}. Sono stati definite macro  \LaTeX per i documenti in modo da facilitare la redazione dei documenti:
  \begin{itemize}
    \item \textbf{Norme di Progetto}: \latex{\textbackslash NdP}
    \item \textbf{Studio di Fattibilità}: \latex{\textbackslash SdF}
    \item \textbf{Analisi dei Requisiti}: \latex{\textbackslash AdR}
    \item \textbf{Piano di Progetto}: \latex{\textbackslash PdP}
    \item \textbf{Piano di Qualifica}: \latex{\textbackslash PdQ}
    \item \textbf{Glossario}: \latex{\textbackslash GlO}
  \end{itemize}
\end{itemize}

\subsubsubsection{Formati}
\begin{itemize}
  \item \textbf{Nome files}: per i nomi dei file deve essere utilizzato il comando \LaTeX: \latex{\textbackslash fileName}
  \item  \textbf{Percorsi}: per i percorsi dei file deve essere utilizzato il comando \LaTeX: \latex{\textbackslash filePath}
  \item \textbf{URLs e Links}:  per gli indirizzi web comleti deve essere utilizzato il comando \LaTeX: \latex{\textbackslash url}, invece per i link deve essere utilizzato il comando \LaTeX: \latex{\textbackslash href}
  \item \textbf{Date e Orari}: le date e gli orari devono seguire la notazione definita dallo standard \emph{ISO 8601}\footnote{\url{http://www.iso.org/iso/iso8601}}: YYYY-MM-DD per le date e HH:MM pe gli orari
  \item \textbf{Quantità e Unità di misura}: le quantità e le unità di misura dovranno aderire allo standard \emph{ISO 80000\footnote{\TODO{link a iso 80000}}}
\end{itemize}

\subsubsubsection{Tabelle}
Ogni tabella presente all’interno dei documenti dev’essere accompagnata da una didascalia, in cui deve comparire un numero identificativo incrementale per la tracciabilità della stessa all’interno del documento.

\subsubsubsection{Immagini e grafici}
Le immagini e i grafici presenti all’interno dei documenti devono essere accompagnate da una didascalia, in cui deve comparire un numero identificativo incrementale per la tracciabilità della stessa all’interno del documento. Le immagini e i grafici devono essere nel formato Scalable Vector Graphics (\gloss{SVG}), così in questo modo si garantisce una maggior qualità dell’immagine in caso di ridimensionamento. Per consentire l’inclusione delle immagini e dei grafici nei documenti, essi dovranno essere convertiti nel formato \gloss{PDF} . Qualora non sia possibile salvare le immagine e i grafici in formato vettoriale è preferito il formato Portable Network Graphics (\gloss{PNG}).

\subsection{Linguaggio}
Nei documenti alcuni termini devono essere interpretati in modo analogo ai termini inglesi descritti in RFC 2119\footnote{\url{http://www.ietf.org/rfc/rfc2119.txt}}:
\begin{itemize}
  \item i termini \emph{deve}, \emph{è richiesto} e sinonimi stretti sono da intendersi con lo stesso significato di \emph{MUST};
  \item i termini \emph{non deve}, \emph{è richiesto che non} e sinonimi stretti sono da intendersi con lo stesso significato di \emph{MUST NOT};
  \item i  termini \emph{dovrebbe}, \emph{si raccomanda}, \emph{è preferibile} e sinonimi stretti sono da intendersi con lo stesso significato di \emph{SHOULD}
  \item i termini \emph{non dovrebbe}, \emph{si raccomanda di non}, \emph{è preferibile che non} e sinonimi stretti sono da intendersi con lo stesso significato di \emph{SHOULD NOT};
  \item i termini \emph{può}, \emph{opzionalmente} e sinonimi stretti sono da intendersi con lo stesso significato di \emph{MAY}.
\end{itemize}

\subsection{Configurazione} %ISO 12207 6.2

\subsection{Verifica} %ISO 12207 6.4

\subsection{Validazione} %ISO 12207 6.5

\subsection{Risoluzione dei Problemi} %ISO 12207 6.8

\newpage

\section{Processi di Organizzazione} %ISO 12207 7.1

Tali processi eseguono funzioni a livello organizzativo aziendale, per supportare altri processi primari, di supporto o organizzativi.
I processi organizzativi aiutano nel definire, controllare e migliorare gli altri processi.

\subsection{Processi di Management}

\subsubsection{Scopo} %ISO 12207 7.1.1

I processi di management contengono le attività generiche e i task, che possono essere impiegati da un qualsiasi gruppo che deve gestire i suoi rispettivi processi. Il \PM{} responsabile per la gestione del prodotto, la gestione del progetto, la gestione dei \textit{\gloss{tasks}} inerenti ai processi come quelli di acquisizione, fornitura, sviluppo, manutenzione o i processi di supporto.

\subsubsection{Descrizione} %ISO 12207 7.1.1

Verranno definite norme riguardanti:
\begin{itemize}
  \item Catalogazione dei ruoli e i loro compiti specifici;
  \item L’identificazione delle modalità di comunicazione interne ed esterne al gruppo;
  \item Definizione di metodologie di Esecuzione e Controllo.
\end{itemize}

\subsubsection{Ruoli di progetto}
%non sono sicuro che se questa sezione sia in qualche modo separata o ricada dentro la pianificazione - 4. Don't Panic

Durante lo sviluppo del progetto vi saranno diversi ruoli che i membri del gruppo andranno a ricoprire. Tali ruoli rappresentano figure aziendali specializzate, indispensabili per il buon esito del progetto. Ciascun componente del gruppo dovrà ricoprire almeno una volta ogni ruolo. Si deve inoltre certificare che non vi siano conflitti di interesse nello svolgimento delle attività di verifica e di approvazione.
Per garantire che la rotazione dei ruoli non provochi conflitti è necessario che le attività di stesura e verifica vengano pianificate dettagliatamente e che i soggetti interessati rispettino i compiti a loro assegnati. Sarà poi compito del \VR{} controllare attentamente il diario delle modifiche di ogni documento per individuare eventuali incongruenze.
Si descrivono ora i diversi ruoli di progetto, con le relative responsabilità e le modalità operative affinché essi possano svolgere i compiti assegnati con l’ausilio dei software scelti per il progetto.

\subsubsubsection{Project Manager}
Il \PM{} rappresenta il progetto, in quanto accentra su di sé le responsabilità di scelta ed approvazione, ed il gruppo, in quanto presenta al committente i risultati del lavoro svolto.
Detiene il potere decisionale, quindi la responsabilità su:
\begin{itemize}
  \item Pianificazione, coordinamento e controllo delle attività;
  \item Gestione e controllo delle risorse;
  \item Analisi e gestione dei rischi;
  \item Approvazione dei documenti;
  \item Approvazione dell'offerta.
\end{itemize}

Di conseguenza, ha il compito di assicurarsi che le attività di verifica vengano svolte sistematicamente seguendo le Norme di Progetto, vengano rispettati i ruoli e le competenze assegnate nel \PdP{}, non vi siano conflitti di interesse tra \textit{\gloss{Redattori}} e
\textit{\gloss{Verificatori}}. Egli è l’unico a poter decidere l’approvazione di un documento e a sancirne la distribuzione.
Ha inoltre l’incarico di gestire la creazione e l’assegnazione dei \textit{\gloss{ticket}} delle macro-fasi e di assegnare ad un membro del gruppo il ruolo di responsabile di quest’ultima.
Redige il \PdP{} e collabora alla stesura del \PdQ{}, in particolare nella sezione relativa alla pianificazione.

\subsubsubsection{Amministratore}
L’\AM{} è responsabile del controllo, dell’efficienza e dell’operatività dell’ambiente di lavoro. Le mansioni di primaria importanza che gli competono sono:
\begin{itemize}
  \item Ricerca di strumenti che possano automatizzare qualsiasi compito che possa essere tolto all’umano;
  \item Risoluzione dei problemi legati alle difficoltà di gestione e controllo dei processi e delle risorse. La risoluzione di tali problemi richiede l’adozione di strumenti adatti;
  \item Controllo delle versioni e delle configurazioni del prodotto;
  \item Gestione dell’archiviazione e del versionamento della documentazione di progetto;
  \item Fornire procedure e strumenti per il monitoraggio e la segnalazione per il controllo qualità.
\end{itemize}
Redige le \NdP{}, dove spiega e norma l’utilizzo degli strumenti, redige la sezione del \PdQ{} dove vengono descritti strumenti e metodi di verifica.

\subsubsubsection{Analista}
L’\AN{} è responsabile delle attività di analisi. Le responsabilità di spicco per tale ruolo sono:
\begin{itemize}
  \item Produrre una specifica di progetto comprensibile, sia per il \gloss{Proponente}, sia per il \textit{\gloss{Committente}} che per il \PG{}, e motivata in ogni suo punto;
  \item Comprendere appieno la natura e la complessità del problema.
\end{itemize}
Redige lo \SdF{}, l’\AdR{} e parte del \PdQ{}.
Partecipa alla redazione del \PdQ{} in quanto conosce l’ambito del progetto ed ha chiari i livelli di qualità richiesta e le procedure da applicare per ottenerla.

\subsubsubsection{Progettista}
Il \PG{} è responsabile delle attività di progettazione. Le responsabilità di tale
ruolo sono:
\begin{itemize}
  \item Produrre una soluzione attuabile, comprensibile e motivata;
  \item Effettuare scelte su aspetti progettuali che applichino al prodotto soluzioni note ed ottimizzate;
  \item Effettuare scelte su aspetti progettuali e tecnologici che rendano il prodotto facilmente manutenibile.
\end{itemize}
Redige la \textit{Specifica Tecnica}, la \textit{Definizione di Prodotto} e le sezioni inerenti le metriche di verifica della programmazione del \PdQ{}.

\subsubsubsection{Programmatore}
Il \PR{} è responsabile delle attività di codifica e delle componenti di ausilio necessarie per l’esecuzione delle prove di verifica e validazione. Le responsabilità di tale ruolo sono:
\begin{itemize}
  \item Implementare rigorosamente le soluzioni descritte dal \PG{}, da cui seguirà quindi la realizzazione del prodotto;
  \item Scrivere codice documentato, versionato, manutenibile e che rispetti gli standard stabiliti per la scrittura del codice;
  \item Implementare i test sul codice scritto, necessari per prove di verifica e validazione.
\end{itemize}
Redige il \textit{Manuale Utente} e produce una abbondante documentazione del codice.

\subsubsubsection{Verificatore}
Il \VR{} è responsabile delle attività di verifica. Ha il compito di effettuare la verifica dei documenti utilizzando gli strumenti e i metodi proposti dal \PdQ{} e attenendosi a quanto descritto nelle \NdP{}. 
Le responsabilità di tale ruolo sono:
\begin{itemize}
  \item Implementare rigorosamente le soluzioni descritte dal \PG{}, da cui seguirà quindi la realizzazione del prodotto;
  \item Assicurare che l’attuazione delle attività sia conforme alle norme stabilite;
  \item Controllare la conformità di ogni stadio del ciclo di vita del prodotto.
\end{itemize}
Redige la sezione del \PdQ{} che illustra l’esito e la completezza delle verifiche
e delle prove effettuate.

\subsubsection{Comunicazione}
%non sono sicuro che se questa sezione sia in qualche modo separata o ricada dentro la pianificazione
\subsubsubsection{Esterna}
Per le comunicazioni esterne è stata creata una casella di posta elettronica:
\subsubsubsection{Interna}

\subsubsection{Riunioni}
%non sono sicuro che se questa sezione sia in qualche modo separata o ricada dentro la pianificazione
\subsubsubsection{Frequenza}

\subsubsubsection{Convocazione riunione}
Il \PM{} ha il compito di convocare le riunioni generali, ovvero le riunioni in cui vengono convocati tutti i membri del gruppo.
Qualora ve ne sia la necessità, qualsiasi componente può richiedere la convocazione di una riunione, ma tale richiesta deve essere inoltrata al \PM{} che a cui spetta il compito di vagliarla. È inoltre possibile e auspicabile che possano essere necessarie riunioni tra specifici membri: ad esempio, nella fase di progettazione può essere utile la collaborazione tra Progettista e Analista, senza richiedere la presenza di persone non direttamente coinvolte, che verranno comunque informate delle decisioni prese tramite verbale.
Il \PM{} deve convocare la riunione con attraverso la creazione di un \textit{\gloss{task}} su \textit{\gloss{Asana}}, assegnato a tutti i membri del gruppo, avente le seguenti caratteristiche:
\begin{itemize}
  \item \textbf{Titolo : } Convocazione riunione n. X, dove X indica il numero crescente di riunioni effettuate;
  \item \textbf{Data : } La data del task dovrà essere impostata alla data prevista per la riunione;
  \item \textbf{Ordine del Giorno : } elenco ordinato delle varie voci da esaminare sotto forma di commento al task.
\end{itemize}
Ogni componente del gruppo deve confermare la propria presenza al \PM{} il prima possibile. 
In caso di mancata risposta o di immotivata assenza il \PM{} ha il dovere di contattare anche telefonicamente colui che non ha fornito una risposta in tempo utile. Il \PM{}, una volta ricevute le risposte e verificata la presenza o assenza dei membri richiesti, può decidere se confermare, annullare o spostare la riunione, per permettere la partecipazione agli eventuali assenti. La conferma, l’annullamento e lo spostamento dell’assemblea devono essere effettuati attraverso modifiche o cancellazione del \textit{\gloss{task}} su \textit{\gloss{Asana}}.
\subsubsubsection{Verbale}
All’apertura della riunione, verificata la presenza dei membri previsti, viene scelto un segretario che avrà il compito di annotare ogni argomento trattato e di redigere il verbale dell’assemblea.
Il segretario deve inoltre controllare che venga seguito l’ordine del giorno in modo da non tralasciare alcun punto.
\paragraph{Riunione interna :}
Il verbale di riunione interna è un documento interno informale che consente di tracciare gli argomenti discussi durante la riunione. Verrà redatto dal segretario della riunione, ruolo scelto di volta in volta a rotazione tra i presenti. Si consiglia di redigere il verbale di una riunione interna in un documento di testo da allegare come commento al \textit{\gloss{task}} creato per la riunione su \textit{\gloss{Asana}}, in modo da rendere sempre disponibile il contenuto dello stesso.

\paragraph{Riunione esterna :}
In caso di riunione con il committente od il proponente, il verbale è un documento ufficiale che può avere valore normativo e quindi deve essere redatto seguendo criteri specifici.
Per agevolare la scrittura di tale documento è stato preparato un template \LaTeX{} che ne definisce la struttura e ne organizza i contenuti. Vi è quindi l’obbligo di seguire il sopraccitato schema per creare e condividere il verbale.
\subsubsection{Pianificazione} %ISO 12207 7.1.2
%task ruoli e responsabilità - Project Libre

\subsubsection{Esecuzione e controllo} %ISO 12207 7.1.3
%come viene monitorata l'esecuzione della pianificazione: DA DEFINIRE (per esempio chiedendo ogni tanto come sta andando il lavoro) - Asana , Git

\newpage

\subsection{Processi di gestione dell'infrastruttura} %ISO 12207 7.2
%asana per i task, github per le issues e bug e risoluzione problemi, projectlibre per la pianificazione con gantt e costi
\subsubsection{Scopo} 
Sono processi per stabilire e mantenere le infrastrutture necessarie a qualsiasi altro processo. Le infrastrutture  possono includere HardWare, SoftWare, strumenti, tecniche, standards, strutture per lo sviluppo o manutenzione.

\subsubsection{Descrizione} 
Verranno definite norme riguardanti:
\begin{itemize}
  \item Definizione e documentazione sulle modalità di \textit{\gloss{ticketing}} ;
  \item Definizione e documentazione sulle modalità di \textit{\gloss{versioning}};
  \item Definizione e documentazione sulla struttura del sistema di \textit{\gloss{repository}}.
\end{itemize}

\subsubsection{Ticketing} 
\subsubsection{Versioning} %Scelta di Git come sistema di Versioning
\subsubsection{Repository} %Scelta di Git come sistema di Repository, Spiegazione albero delle directory 

\newpage


\end{document}
